\subsubsection{常量定义头文件}
在include目录下的fat.h文件中定义所有用到的常量、结构体如下:
\begin{ccode}{常量与结构体定义}
#include<stdio.h>
#include<string.h>
#include<stdlib.h>
#include<stdbool.h>

// 定义所用的常量
#define DISK_SIZE 2560//磁盘大小
#define BLOCK_SIZE 64 //磁盘块大小
#define BLOCK_NUM  (DISK_SIZE / BLOCK_SIZE)
#define MAX_ENTRY 32
#define MAX_OPENFILE 32
#define MAX_FD 16
#define DIR_NAME_LENGTH 16
// 定义模式
#define MODE_READ 1
#define MODE_WRITE 2
#define MODE_RDWR 3
#define MODE_WAPPEND 4  //写追加
// 全局数组声明
extern char Disk[DISK_SIZE];
extern int FAT[BLOCK_NUM]; // FAT表,个数即为块数。
extern int current_dir;//当前目录索引


typedef struct{
    char name[DIR_NAME_LENGTH]; // 目录长度
    int size;
    size_t start_block;
    int ac;//是否使用,0表示未使用,1表示已使用
    int type;//0表示文件,1表示目录
    int parent;
}DirEntry;

typedef struct{
    int used;//0空闲,1已经使用
    int dir_index;//文件下标
    int offset;//读写偏移
    int mode;//打开模式
}FileDescriptor;

extern DirEntry Directory[MAX_ENTRY];
extern FileDescriptor FDTable[MAX_FD]; 
\end{ccode}
上述代码中定义了磁盘大小、磁盘块大小、磁盘块数(磁盘大小/磁盘块大小)以及最大目录项数、最大打开文件数、最大文件描述符数和最大文件名长度。
定义了目录项结构体DirEntry,包含属性目录名、目录大小、起始磁盘块号、是否使用、目录类型(0表示文件,1表示目录)和父目录索引。
定义了文件描述符结构体FileDescriptor,包含属性是否使用、文件下标、读写偏移以及打开模式(对应只读r、只写w和读写rw)。还声明了全局变量
磁盘数组Disk用以描述磁盘;FAT数组表示磁盘块的链接和分配关系;current\_dir表示当前目录位置索引;Directory数组表示目录项;FDTable数组表示文件描述符。